\section{Conditional expectation from a faithful fixed point}

This section follows~\cite[Section~6.4]{Wolf2012Quantum}; the conditional
expectation is built from the fixed-point algebra of~\cite[Theorem~6.14]{Wolf2012Quantum}
via~\cite[(1.40)]{Wolf2012Quantum}.
The scalar conditional expectation $E_\sigma$ is the special case where
the fixed-point $*$-subalgebra is the scalar algebra $\C\cdot\Id$
(the primitive-channel case). The general irreducible case with
nontrivial period requires the Wedderburn block decomposition
of~\cite[Theorem~6.14]{Wolf2012Quantum}.
% The nontrivial-period conditional expectation is outside the present
% formal scope of this section.

\begin{definition}[Conditional expectation]\label{def:is_conditional_expectation}
    \lean{Kraus.IsConditionalExpectation}
    \leanok
    Given a $*$-subalgebra $S \subseteq \MN{D}$, a linear map
    $P : \MN{D} \to \MN{D}$ is a conditional expectation onto $S$ if it is
    idempotent, unital, has range contained in $S$, and satisfies
    $P(X)=X$ for all $X \in S$.
\end{definition}

\begin{definition}[Scalar conditional expectation]\label{def:scalar_conditional_expectation}
    \lean{Kraus.scalarConditionalExpectation}
    \leanok
    \uses{def:is_conditional_expectation}
    For $\sigma \ge 0$ with $\tr(\sigma)\neq 0$, define
    \[
        E_\sigma(X) := \frac{\tr(\sigma X)}{\tr(\sigma)} \Id.
    \]
\end{definition}

\begin{lemma}[Evaluation formula]\label{lem:scalar_conditional_expectation_apply}
    \lean{Kraus.scalarConditionalExpectation_apply}
    \leanok
    \uses{def:scalar_conditional_expectation}
    $E_\sigma(X)=\dfrac{\tr(\sigma X)}{\tr(\sigma)} \Id$.
\end{lemma}

\begin{lemma}[Unitality]\label{thm:scalar_conditional_expectation_unital}
    \lean{Kraus.scalarConditionalExpectation_unital}
    \leanok
    \uses{def:scalar_conditional_expectation}
    $E_\sigma(\Id)=\Id$.
\end{lemma}

\begin{lemma}[Idempotence]\label{thm:scalar_conditional_expectation_idempotent}
    \lean{Kraus.scalarConditionalExpectation_idempotent}
    \leanok
    \uses{def:scalar_conditional_expectation}
    $E_\sigma(E_\sigma(X)) = E_\sigma(X)$ for all $X$.
\end{lemma}

\begin{lemma}[Range is scalar]\label{lem:scalar_conditional_expectation_range_scalar}
    \lean{Kraus.scalarConditionalExpectation_range_scalar}
    \leanok
    \uses{def:scalar_conditional_expectation}
    For every $X$, there exists $c\in\C$ with $E_\sigma(X)=c\,\Id$.
\end{lemma}

\begin{lemma}[Fixes scalar operators]\label{lem:scalar_conditional_expectation_fixes_scalar}
    \lean{Kraus.scalarConditionalExpectation_fixes_scalar}
    \leanok
    \uses{def:scalar_conditional_expectation}
    For $c\in\C$, $E_\sigma(c\,\Id)=c\,\Id$.
\end{lemma}

\begin{lemma}[Absorbs the adjoint map]\label{thm:scalar_conditional_expectation_absorbs_adjointMap}
    \lean{Kraus.scalarConditionalExpectation_absorbs_adjointMap}
    \leanok
    \uses{def:scalar_conditional_expectation}
    If $T(\sigma)=\sigma$, then
    $E_\sigma\circ T^* = E_\sigma$.
\end{lemma}

\begin{lemma}[Adjoint map absorbs scalar projection]
    \label{thm:adjointMap_absorbs_scalarConditionalExpectation}
    \lean{Kraus.adjointMap_absorbs_scalarConditionalExpectation}
    \leanok
    \uses{def:scalar_conditional_expectation, def:tp_kraus}
    If $T$ is trace-preserving, then
    $T^*\circ E_\sigma = E_\sigma$.
\end{lemma}

\begin{lemma}[Image lies in adjoint fixed points]
    \label{lem:scalar_conditional_expectation_mem_adjointFixedPoints}
    \lean{Kraus.scalarConditionalExpectation_mem_adjointFixedPoints}
    \leanok
    \uses{thm:adjointMap_absorbs_scalarConditionalExpectation}
    For all $X$, $E_\sigma(X)$ is fixed by $T^*$.
\end{lemma}

\begin{theorem}[Scalar map is a conditional expectation]
    \label{thm:scalar_conditional_expectation_isConditionalExpectation}
    \lean{Kraus.scalarConditionalExpectation_isConditionalExpectation}
    \leanok
    \uses{def:is_conditional_expectation, def:scalar_conditional_expectation,
        def:tp_kraus,
        thm:scalar_conditional_expectation_unital,
        thm:scalar_conditional_expectation_idempotent,
        lem:scalar_conditional_expectation_mem_adjointFixedPoints,
        lem:scalar_conditional_expectation_fixes_scalar}
    Let $K$ be trace-preserving, let $\rho>0$ be a positive-definite fixed
    point of the associated channel $T(\rho)=\sum_i K_i\rho K_i^\dagger$ (so $T(\rho)=\rho$), and
    assume every adjoint fixed point of $T^*$ is a scalar multiple of $\Id$.
    Then $E_\rho$ is a conditional expectation onto the adjoint fixed-point
    $*$-subalgebra.
\end{theorem}

\begin{lemma}[Complete positivity of the scalar conditional expectation]
    \label{thm:scalar_conditional_expectation_isCPMap}
    \lean{Kraus.scalarConditionalExpectation_isCPMap}
    \leanok
    \uses{def:scalar_conditional_expectation}
    For any positive semidefinite $\sigma$, the scalar conditional
    expectation $E_\sigma$ is a completely positive map.
\end{lemma}
\begin{proof}\leanok
    Write $S=\sqrt{\sigma}$ and define the Kraus family
    $K_{j,i}=\ketbra{j}{i}\,S$ indexed by $(j,i)\in\{1,\ldots,D\}^2$.
    Using $\ketbra{j}{i}\,M\,\ketbra{i}{j} = M_{ii} \ketbra{j}{j}$ and
    summing first over $i$ and then over $j$,
    \[
        \sum_{j,i} K_{j,i}\,X\,K_{j,i}^\dagger
        = \sum_j \ketbra{j}{j} \tr(SXS)
        = \tr(\sigma X) \Id,
    \]
    where we used cyclicity of the trace and $S^2=\sigma$. Hence
    $E_\sigma = (\tr(\sigma))^{-1}\cdot(X\mapsto\sum_{j,i}K_{j,i}XK_{j,i}^\dagger)$
    is the Kraus map scaled by the non-negative real number
    $(\tr(\sigma))^{-1}$, and therefore completely positive.
\end{proof}

\begin{lemma}[$\sigma$-trace preservation of the scalar conditional expectation]
    \label{thm:scalar_conditional_expectation_preserves_weighted_trace}
    \lean{Kraus.scalarConditionalExpectation_preserves_weighted_trace}
    \leanok
    \uses{def:scalar_conditional_expectation}
    If $\tr(\sigma)\neq 0$, then for every $X$
    \[
        \tr(\sigma\,E_\sigma(X))
        = \tr(\sigma X).
    \]
\end{lemma}
\begin{proof}\leanok
    Direct computation:
    $\tr(\sigma\cdot\tfrac{\tr(\sigma X)}{\tr(\sigma)}\Id)
        = \tfrac{\tr(\sigma X)}{\tr(\sigma)}\tr(\sigma)
        = \tr(\sigma X)$.
\end{proof}

\section{Stationary support}

This section develops the support-projection theory
behind stationary states, following~\cite[Section~6.4, Propositions~6.10--6.11]{Wolf2012Quantum}.

\begin{lemma}[Lower-triangular Kraus condition implies corner invariance]%
    \label{lem:lower_zero_invariance}
    \lean{MPSTensor.lowerZero_implies_invariance}
    \leanok
    Let $A$ be an MPS tensor and let $P$ be an orthogonal projection.
    If, for every physical index~$i$,
    \[
        (\Id - P) A^i P = 0
    \]
    then
    \[
        P \E_A(P X P) P = \E_A(P X P)
    \]
    for all $X \in \MN{D}$.
\end{lemma}

\begin{proof}\leanok
    Expand the transfer map and use $(\Id - P) A^i P = 0$ to reduce
    each Kraus term to the $P$-corner.
\end{proof}

\begin{lemma}[Support projection invariance]%
    \label{thm:support_proj_fixed}
    \lean{Channel.support_proj_fixed}
    \leanok
    \uses{def:channel}
    Let $E$ be a quantum channel and $\rho \ge 0$ a positive
    semidefinite fixed point ($E(\rho) = \rho$).
    Let $P$ denote the support projection of~$\rho$.
    Then
    \[
        P E(P X P) P = E(P X P)
    \]
    for all $X \in \MN{D}$.
\end{lemma}

\begin{proof}
    \leanok
    \uses{lem:lower_zero_invariance}
    Write $E$ in Kraus form $E(X) = \sum_i K_i X K_i^\dagger$.
    From $E(\rho) = \rho$ and $\rho \ge 0$ one deduces
    $(\Id - P) K_i P = 0$ for all~$i$.
    Apply Lemma~\ref{lem:lower_zero_invariance}.
\end{proof}

\begin{definition}[Stationary state]\label{def:stationaryState}
    \lean{Channel.stationaryState}
    \leanok
    \uses{def:channel, def:irreducible_cp}
    For an irreducible channel $E$ on $\MN{D}$ (with $D \ge 1$),
    the \emph{stationary state} is the unique density-matrix fixed
    point $\rho_\infty$ satisfying $E(\rho_\infty) = \rho_\infty$,
    $\rho_\infty > 0$, and $\tr(\rho_\infty) = 1$.
    Uniqueness follows from the quantum Perron--Frobenius theorem
    (Theorem~\ref{thm:qpf}).
\end{definition}

\begin{definition}[Stationary support]\label{def:stationarySupport}
    \lean{Channel.stationarySupport}
    \leanok
    \uses{def:stationaryState}
    The \emph{stationary support} of an irreducible channel~$E$
    is the support projection of the stationary state
    (Definition~\ref{def:stationaryState}).
\end{definition}

\begin{theorem}[Stationary support is full]%
    \label{thm:stationarySupport_eq_one}
    \lean{Channel.stationarySupport_eq_one}
    \leanok
    \uses{def:stationarySupport, def:irreducible_cp}
    For an irreducible channel~$E$, the stationary support equals
    the identity: $\mathrm{supp}(\rho_\infty) = \Id$.
\end{theorem}

\begin{proof}
    \leanok
    \uses{thm:support_proj_fixed}
    The support projection $P$ of $\rho_\infty$ is an orthogonal
    projection satisfying $P E(PXP) P = E(PXP)$ for all~$X$
    (Lemma~\ref{thm:support_proj_fixed}).
    By irreducibility, $P = 0$ or $P = \Id$.
    Since $\rho_\infty \neq 0$, we have $P \neq 0$, so $P = \Id$.
\end{proof}

\begin{remark}[Stationary-support formulations]%
    \label{rem:stationary_support_todo}
    The non-trivial formulations needed here are the following results from
    Wolf:
    \begin{itemize}
        \item the converse direction of~\cite[Proposition~6.10]{Wolf2012Quantum}:
              if the support projection of every positive semidefinite fixed point
              is full, then the channel is irreducible;
        \item \cite[Proposition~6.11]{Wolf2012Quantum}: minimality of the
              stationary support among invariant projections, without assuming
              irreducibility.
    \end{itemize}
    These statements are not used elsewhere in this chapter.
\end{remark}

\subsection{Faithful compression onto the support sector}

This subsection establishes the results on the corner algebra,
$*$-structure, and compression of a PSD fixed point onto its support sector
needed for~\cite[Corollary~6.6]{Wolf2012Quantum}: the corner carries a
$\C$-algebra and $*$-structure, the two standard presentations of the corner
are identified, and the compression of a PSD fixed point onto its support
sector is positive definite.

\begin{lemma}[Support projection annihilates the kernel]%
    \label{lem:supportProj_mulVec_eq_zero_of_mulVec_eq_zero}
    \lean{MPSTensor.supportProj_mulVec_eq_zero_of_mulVec_eq_zero}
    \leanok
    Let $\rho \succeq 0$ on $\C^D$ with support projection
    $P = \mathrm{supp}(\rho)$. For every $v \in \C^D$, if $\rho\, v = 0$,
    then $P\, v = 0$.
\end{lemma}
\begin{proof}\leanok
    Diagonalize $\rho = U\,\diag(\lambda)\,U^\dagger$ with
    $\lambda_j \ge 0$. Writing $w := U^\dagger v$, the hypothesis
    $\rho v = 0$ gives $\lambda_j\, w_j = 0$ for every $j$; hence $w_j = 0$
    whenever $\lambda_j > 0$. The support projection acts as
    $\diag(\chi_{\lambda_j > 0})$ in the eigenbasis, so
    $P\,v = U\,\diag(\chi_{\lambda_j > 0})\,w = 0$.
\end{proof}

\begin{definition}[Corner submodule / idempotent-corner identification]%
    \label{def:cornerSubmoduleCornerEquiv}
    \lean{MatrixCorner.cornerSubmoduleCornerEquiv}
    \leanok
    For an idempotent $P \in \MN{D}$, the subspace
    $\{X \mid P X P = X\}$ and the corner
    $\{P\,X\,P \mid X \in \MN{D}\}$ are identified as $\C$-linear spaces by
    the identity on underlying matrices.
\end{definition}

\begin{remark}[$\C$-algebra and $*$-structure on the corner]%
    \label{rem:corner_algebra_star_instances}
    \lean{MatrixCorner.instAlgebraComplexCorner}
    \lean{MatrixCorner.cornerStar}
    \lean{MatrixCorner.cornerStarRing}
    \lean{MatrixCorner.cornerStarModuleComplex}
    \leanok
    The corner algebra associated to an idempotent already carries a ring
    structure with unit~$P$. In the matrix case it also has
    $\C$-module and $\C$-algebra structures and, under the additional
    assumption $P^\dagger = P$, star, star ring, and star module
    structures over~$\C$.
\end{remark}

\begin{theorem}[Faithful compression onto the support]%
    \label{thm:compression_on_support_posDef}
    \lean{Matrix.PosSemidef.compression_on_support_posDef}
    \leanok
    \uses{lem:supportProj_mulVec_eq_zero_of_mulVec_eq_zero}
    Let $\rho \succeq 0$ on $\C^D$ with support projection
    $P = \mathrm{supp}(\rho)$, and let
    $V : \C^D \to \C^k$ be a compression isometry satisfying
    $V\,V^\dagger = \Id_k$ and $V^\dagger\,V = P$. Then the compression
    \[
        V\,\rho\,V^\dagger \in \MN{k}
    \]
    is positive definite.
\end{theorem}
\begin{proof}\leanok
    \uses{lem:supportProj_mulVec_eq_zero_of_mulVec_eq_zero}
    Hermiticity follows from Hermiticity of $\rho$ and the product-star
    identity. For strict positivity, let $w \neq 0$, set
    $u := V^\dagger w$, and note that $V V^\dagger = \Id_k$ forces
    $u \neq 0$, while $V^\dagger V = P$ forces $P u = u$. The quadratic
    form computes as
    $w^\dagger (V\rho V^\dagger) w = u^\dagger \rho u$,
    which is non-negative by positive semidefiniteness. If it vanishes, then
    $\rho u = 0$, and by
    Lemma~\ref{lem:supportProj_mulVec_eq_zero_of_mulVec_eq_zero} also
    $P u = 0$, contradicting $P u = u \neq 0$.
\end{proof}

\begin{theorem}[Corner-restricted fixed points form a $*$-algebra]%
    \label{thm:cornerFixedPointsStarSubalgebra}
    \lean{Kraus.cornerFixedPointsStarSubalgebra}
    \leanok
    \uses{thm:compression_on_support_posDef, rem:corner_algebra_star_instances,
        thm:adjointFixedPointsStarSubalgebra, thm:support_proj_fixed}
    Let $T^*(Y) = \sum_i K_i^\dagger Y K_i$ be a trace-preserving Schwarz map,
    let $\rho \succeq 0$ satisfy $T(\rho) = \rho$ for the Schr\"odinger map
    $T(X) = \sum_i K_i X K_i^\dagger$, and let $Q = \mathrm{supp}(\rho)$ be the
    support projection of~$\rho$. Then
    \[
        \{Y \in Q\,\MN{D}\,Q \mid Q\,T^*(Y)\,Q = Y\}
    \]
    is a $*$-subalgebra of the corner algebra $Q\,\MN{D}\,Q$.
\end{theorem}
\begin{proof}
    \leanok
    \uses{thm:compression_on_support_posDef, thm:adjointFixedPointsStarSubalgebra,
        thm:support_proj_fixed}
    The stated set is the fixed-point set of the compressed map $\widetilde{T}^*$
    on the support sector. Since $\rho$ is a fixed point of $T$, the support
    projection satisfies $(\Id - Q) K_i Q = 0$
    (Theorem~\ref{thm:support_proj_fixed}), so the compressed family
    $Q K_i Q$ is supported on the corner and is trace-preserving there,
    $\sum_i (Q K_i Q)^\dagger (Q K_i Q) = Q$. Compress this family to the
    support sector: the compressed map is unital, and the compression of $\rho$
    onto the sector is positive definite
    (Theorem~\ref{thm:compression_on_support_posDef}) and is a fixed point of
    the compressed Schr\"odinger map. The structure theorem for fixed points of
    unital Schwarz maps with a positive-definite fixed point
    (Theorem~\ref{thm:adjointFixedPointsStarSubalgebra}) makes the compressed
    fixed points a $*$-algebra. The compression isomorphism is multiplicative
    and star-preserving, intertwines the compressed adjoint map with the
    corner-restricted map, and transports this $*$-algebra structure back to the
    corner.
\end{proof}

\section{Wedderburn decomposition of the fixed-point algebra}

This section proves the structure theorem for the fixed-point
algebra of a trace-preserving Kraus map, following~\cite[Theorems~6.12--6.14]{Wolf2012Quantum}. The proof uses the
Jacobson radical characterization of semisimplicity together with the
Wedderburn--Artin structure theorem for finite-dimensional semisimple
algebras over an algebraically closed field.

The algebraic core is the semisimplicity of an arbitrary $*$-subalgebra of a
complex matrix algebra, from which the abstract Wedderburn--Artin decomposition
follows. The fixed-point algebra of a trace-preserving Kraus map is then the
special case used in this section.

\begin{lemma}[A $*$-subalgebra of a matrix algebra is semisimple]%
    \label{lem:starSubalgebra_isSemisimpleRing}
    \lean{StarSubalgebra.isSemisimpleRing}
    \leanok
    Every $*$-subalgebra of $\MN{D}$ is a semisimple ring.
\end{lemma}

\begin{proof}
    \leanok
    The subalgebra is finite-dimensional, hence Artinian, so its Jacobson
    radical $J$ is nilpotent. If $x \in J$ then $x^*x \in J$, since the radical
    is a two-sided ideal. In the matrix algebra $x^*x$ is positive
    semidefinite, and a positive semidefinite nilpotent matrix vanishes: its
    trace is a nilpotent complex number, hence zero, and the trace of $x^*x$
    vanishes only when $x = 0$. Thus $J = 0$ and the subalgebra is semisimple.
\end{proof}

\begin{theorem}[Abstract Wedderburn--Artin decomposition for $*$-subalgebras]%
    \label{thm:starSubalgebra_wedderburnArtin}
    \lean{StarSubalgebra.exists_algEquiv_pi_matrix}
    \leanok
    \uses{lem:starSubalgebra_isSemisimpleRing}
    There exist $n \in \N$ and dimensions $d_1,\ldots,d_n \ge 1$ such that every
    $*$-subalgebra of $\MN{D}$ is $\C$-algebra isomorphic to
    \[
        \prod_{k=1}^{n} \MN{d_k}.
    \]
\end{theorem}

\begin{proof}
    \leanok
    \uses{lem:starSubalgebra_isSemisimpleRing}
    Combine semisimplicity (Lemma~\ref{lem:starSubalgebra_isSemisimpleRing})
    with the Wedderburn--Artin structure theorem for finite-dimensional
    semisimple algebras over the algebraically closed field $\C$.
\end{proof}

The abstract decomposition records only the ring structure. To realize it by a
single unitary change of basis, as in~\cite[Theorem~6.14]{Wolf2012Quantum}, one
uses the spatial feature of a $*$-subalgebra: it reduces orthogonally. Whenever a
subspace is invariant, so is its orthogonal complement, so the representation
space splits into mutually orthogonal invariant pieces rather than merely into a
direct sum of them.

\begin{lemma}[Orthogonal reducibility of a $*$-subalgebra]%
    \label{lem:starSubalgebra_orthogonal_reducibility}
    \lean{StarSubalgebra.orthogonal_mem_invtSubmodule}
    \leanok
    Let $S$ be a $*$-subalgebra of $\MN{D}$, and let $W \subseteq \C^{D}$ be a
    subspace invariant under every member of $S$, meaning $A\,W \subseteq W$ for
    all $A \in S$. Then the orthogonal complement $W^{\perp}$ is also invariant
    under every member of $S$: for each $A \in S$,
    \[
        A\, W^{\perp} \subseteq W^{\perp}.
    \]
\end{lemma}

\begin{proof}
    \leanok
    Fix $A \in S$. Since $S$ is closed under the adjoint, $A^{\dagger} \in S$,
    and therefore $A^{\dagger} W \subseteq W$. For $v \in W^{\perp}$ and
    $w \in W$,
    \[
        \langle w, A v \rangle = \langle A^{\dagger} w, v \rangle = 0,
    \]
    because $A^{\dagger} w \in W$ while $v$ is orthogonal to $W$. Hence
    $A v \in W^{\perp}$.
\end{proof}

\begin{lemma}[Invariant subspaces split off orthogonally]%
    \label{lem:starSubalgebra_isCompl_orthogonal}
    \lean{StarSubalgebra.isCompl_orthogonal_of_invariant}
    \leanok
    \uses{lem:starSubalgebra_orthogonal_reducibility}
    With $S$ and $W$ as in
    Lemma~\ref{lem:starSubalgebra_orthogonal_reducibility}, the representation
    space decomposes as the orthogonal direct sum
    \[
        \C^{D} = W \oplus W^{\perp},
    \]
    in which the orthogonal complement $W^{\perp}$ is invariant under every
    member of $S$.
\end{lemma}

\begin{proof}
    \leanok
    \uses{lem:starSubalgebra_orthogonal_reducibility}
    In a finite-dimensional inner product space the orthogonal complement of a
    subspace is a vector-space complement, which gives the orthogonal direct
    sum. Invariance of $W^{\perp}$ is
    Lemma~\ref{lem:starSubalgebra_orthogonal_reducibility}.
\end{proof}

Orthogonal reducibility is the inductive step behind a full orthogonal
decomposition. A subspace with no proper nonzero invariant subspace is
irreducible; a subspace that has one splits off its orthogonal complement, and
both pieces are again invariant and strictly smaller. Iterating on dimension
exhausts the space into irreducible pieces, mutually orthogonal because at each
step the two pieces are.

\begin{definition}[Irreducible invariant subspace]%
    \label{def:starSubalgebra_irreducibleSubspace}
    \lean{StarSubalgebra.IsIrreducibleSubspace}
    \leanok
    Let $S$ be a $*$-subalgebra of $\MN{D}$. A subspace $W \subseteq \C^{D}$ is
    irreducible under $S$ when it is nonzero, invariant under every member of $S$,
    and its only invariant subspaces are $\{0\}$ and $W$ itself.
\end{definition}

\begin{lemma}[Orthogonal decomposition of an invariant subspace]%
    \label{lem:starSubalgebra_orthogonal_irreducible_invariant}
    \lean{StarSubalgebra.exists_orthogonal_irreducible_of_invariant}
    \leanok
    \uses{lem:starSubalgebra_isCompl_orthogonal, def:starSubalgebra_irreducibleSubspace}
    Let $S$ be a $*$-subalgebra of $\MN{D}$ and let $W \subseteq \C^{D}$ be invariant
    under every member of $S$. Then $W$ is the sum of a finite family of pairwise
    orthogonal subspaces contained in $W$, each irreducible under $S$.
\end{lemma}

\begin{proof}
    \leanok
    \uses{lem:starSubalgebra_isCompl_orthogonal}
    Induct on the dimension of $W$. If $W$ is zero the empty family works, and if
    $W$ is irreducible the family $\{W\}$ works. Otherwise $W$ has a proper nonzero
    invariant subspace $U$. By
    Lemma~\ref{lem:starSubalgebra_isCompl_orthogonal} the orthogonal complement of
    $U$ inside $W$ is invariant, and both $U$ and that complement are strictly
    smaller than $W$ and span it. Their inductive decompositions consist of
    subspaces lying inside $U$ and inside its complement respectively, hence are
    orthogonal to one another; the two families together form the required family.
\end{proof}

\begin{theorem}[Orthogonal irreducible decomposition]%
    \label{thm:starSubalgebra_orthogonal_irreducible_decomposition}
    \lean{StarSubalgebra.exists_orthogonal_irreducible_decomposition}
    \leanok
    \uses{lem:starSubalgebra_orthogonal_irreducible_invariant}
    Let $S$ be a $*$-subalgebra of $\MN{D}$. Then $\C^{D}$ is the sum of a finite
    family of pairwise orthogonal subspaces, each irreducible under $S$; that is,
    $\C^{D}$ is an orthogonal direct sum of irreducible invariant subspaces.
\end{theorem}

\begin{proof}
    \leanok
    \uses{lem:starSubalgebra_orthogonal_irreducible_invariant}
    Apply Lemma~\ref{lem:starSubalgebra_orthogonal_irreducible_invariant} to the
    whole space $W = \C^{D}$, which is invariant under every member of $S$.
\end{proof}

The irreducibility notion just used is simplicity in the language of modules. The
members of $S$ act on $\C^{D}$ as the operators their matrices represent, making
$\C^{D}$ a module over the algebra $S$.

\begin{lemma}[Irreducible invariant subspaces are the simple submodules]%
    \label{lem:starSubalgebra_irreducible_iff_simple}
    \lean{StarSubalgebra.isInvariantSubspace_iff_exists_restrictScalars,
      StarSubalgebra.isIrreducibleSubspace_restrictScalars_iff}
    \leanok
    \uses{def:starSubalgebra_irreducibleSubspace}
    Let $S$ be a $*$-subalgebra of $\MN{D}$, acting on $\C^{D}$. A complex
    subspace $W \subseteq \C^{D}$ is invariant under every member of $S$ if and
    only if $W$ is the underlying complex subspace of a submodule of $\C^{D}$
    over $S$. Under this identification an invariant subspace is irreducible
    under $S$ precisely when, as a module over $S$, it is simple: nonzero, with
    no submodule other than $\{0\}$ and itself.
\end{lemma}

\begin{proof}
    \leanok
    \uses{def:starSubalgebra_irreducibleSubspace}
    The underlying set of a submodule of $\C^{D}$ over $S$ is closed under
    addition and under the action of every member of $S$; since the complex
    scalars lie in $S$, it is a complex subspace, and it is invariant.
    Conversely an invariant complex subspace $W$ is closed under addition and
    under the action of every member of $S$, so $W$ underlies a submodule $q$
    of $\C^{D}$ over $S$ with the same members. Under this correspondence the
    submodules below $q$ are exactly the invariant subspaces below $W$:
    \[
        \{\, q' \subseteq q : q' \text{ a submodule over } S \,\}
        =
        \{\, W' \subseteq W : W' \text{ invariant under } S \,\}.
    \]
    Hence $q$ has a submodule other than $\{0\}$ and $q$ exactly when $W$ has
    an invariant subspace other than $\{0\}$ and $W$: simplicity of $q$ as a
    module over $S$ is irreducibility of $W$ under $S$.
\end{proof}

\begin{lemma}[Euclidean space is a semisimple module over the subalgebra]%
    \label{lem:starSubalgebra_isSemisimpleModule}
    \lean{StarSubalgebra.isSemisimpleModule_euclideanSpace}
    \leanok
    \uses{lem:starSubalgebra_isSemisimpleRing}
    Let $S$ be a $*$-subalgebra of $\MN{D}$. As a module over $S$, the space
    $\C^{D}$ is semisimple: every submodule admits a complementary submodule.
\end{lemma}

\begin{proof}
    \leanok
    \uses{lem:starSubalgebra_isSemisimpleRing}
    The ring $S$ is semisimple (Lemma~\ref{lem:starSubalgebra_isSemisimpleRing}),
    and every module over a semisimple ring is semisimple.
\end{proof}

The next step towards the structure theorem groups the irreducible pieces by
isomorphism type. The grouping is governed by Schur's lemma: a map commuting with
$S$ between two pieces is either zero or an isomorphism, and over the complex
numbers the commutant of $S$ on a single piece is the scalars.

\begin{lemma}[Schur's lemma: an intertwiner of irreducible pieces is zero or an isomorphism]%
    \label{lem:starSubalgebra_intertwiner_irreducible}
    \lean{StarSubalgebra.intertwiner_zero_or_isomorphism}
    \leanok
    \uses{def:starSubalgebra_irreducibleSubspace}
    Let $S$ be a $*$-subalgebra of $\MN{D}$ and let $W, W' \subseteq \C^{D}$ be irreducible
    under $S$. Let $f \colon \C^{D} \to \C^{D}$ be a linear operator that carries $W$ into
    $W'$ and commutes on $W$ with every member of $S$, that is,
    \[
        f(A x) = A\,(f x)
    \]
    for every $A \in S$ and every $x \in W$. Then either $f$ is zero on $W$, or $f$ is
    injective on $W$ and maps $W$ onto $W'$.
\end{lemma}

\begin{proof}
    \leanok
    \uses{def:starSubalgebra_irreducibleSubspace}
    The vectors of $W$ annihilated by $f$ form a subspace of $W$ that is invariant under
    $S$, because $f$ commutes with $S$ on $W$. By irreducibility of $W$ this kernel is
    either all of $W$, in which case $f$ is zero on $W$, or trivial, in which case $f$ is
    injective on $W$. Likewise the image of $W$ under $f$ is a subspace of $W'$ invariant
    under $S$, so by irreducibility of $W'$ it is either trivial, again forcing $f$ to be
    zero on $W$, or all of $W'$. Hence $f$ is zero on $W$, or it is injective on $W$ with
    image $W'$.
\end{proof}

\begin{lemma}[Scalar commutant on an irreducible piece]%
    \label{lem:starSubalgebra_scalar_commutant_irreducible}
    \lean{StarSubalgebra.exists_eq_smul_of_commute_of_irreducible}
    \leanok
    \uses{def:starSubalgebra_irreducibleSubspace}
    Let $S$ be a $*$-subalgebra of $\MN{D}$ and let $W \subseteq \C^{D}$ be irreducible
    under $S$. Let $f \colon \C^{D} \to \C^{D}$ be a linear operator that maps $W$ into
    itself and commutes on $W$ with every member of $S$, that is,
    \[
        f(A x) = A\,(f x)
    \]
    for every $A \in S$ and every $x \in W$. Then $f$ acts on $W$ as a scalar: there is a
    complex number $c$ with
    \[
        f x = c\,x
    \]
    for every $x \in W$.
\end{lemma}

\begin{proof}
    \leanok
    \uses{def:starSubalgebra_irreducibleSubspace}
    Restricting $f$ to $W$ gives an operator on a nonzero finite-dimensional complex
    vector space, which over the complex numbers has an eigenvalue $c$ together with an
    eigenvector inside $W$. The set of vectors of $W$ on which $f$ acts as multiplication
    by $c$ is a subspace contained in $W$; it is nonzero because it contains that
    eigenvector, and it is invariant under $S$ because $f$ commutes with every member of
    $S$ on $W$ while the scalar $c$ commutes with everything. By irreducibility this
    subspace is all of $W$, so $f$ acts as multiplication by $c$ throughout $W$.
\end{proof}

Schur's lemma turns the qualitative notion of two pieces ``looking the same'' into a
usable relation. Call a linear operator on $\C^{D}$ an intertwiner from $W$ into $W'$ when
it carries $W$ into $W'$ and commutes on $W$ with every member of $S$. Two pieces are of
the same type when some intertwiner between them does not vanish on the source; on
irreducible pieces such an intertwiner is automatically an isomorphism.

\begin{definition}[Intertwiner between subspaces]%
    \label{def:starSubalgebra_intertwiner}
    \lean{StarSubalgebra.IsIntertwiner}
    \leanok
    Let $S$ be a $*$-subalgebra of $\MN{D}$ and let $W, W' \subseteq \C^{D}$ be subspaces. A
    linear operator $f \colon \C^{D} \to \C^{D}$ is an intertwiner from $W$ into
    $W'$ when it carries $W$ into $W'$ and, for every $A \in S$ and every $x \in W$,
    \[
        f(A x) = A\,(f x).
    \]
\end{definition}

An intertwiner whose source is irreducible scales the inner product by one nonnegative real
factor; when the intertwiner does not vanish on the source, that factor is positive, so after
rescaling it is an isometry. This is the metric ingredient that turns each same-type
isomorphism into a partial isometry.

\begin{lemma}[Intertwiner scales the inner product]%
    \label{lem:starSubalgebra_intertwiner_inner_scale}
    \lean{StarSubalgebra.exists_real_smul_inner_of_intertwiner}
    \leanok
    \uses{def:starSubalgebra_intertwiner, def:starSubalgebra_irreducibleSubspace,
        lem:starSubalgebra_scalar_commutant_irreducible}
    Let $S$ be a $*$-subalgebra of $\MN{D}$, let $W \subseteq \C^{D}$ be irreducible under $S$,
    and let $f$ be an intertwiner from $W$ into a subspace $W'$. Then there is a
    real number $c \ge 0$ such that for all $x, y \in W$,
    \[
        \langle f(x), f(y) \rangle = c\,\langle x, y \rangle .
    \]
\end{lemma}

\begin{proof}
    \leanok
    \uses{def:starSubalgebra_intertwiner, def:starSubalgebra_irreducibleSubspace,
        lem:starSubalgebra_scalar_commutant_irreducible}
    Following $f$ with its adjoint and the orthogonal projection $P_{W}$ onto $W$ gives an
    operator $g = P_{W} f^{*} f$ that maps $W$ into itself, since the projection lands in $W$.
    On $W$ it commutes with every member of $S$: for $A \in S$ the adjoint $A^{*}$ is again a
    member, and combining the intertwining identity of $f$ with the pairing of the projected
    adjoint against $W$ gives $g(A x) = A\,(g x)$ for $x \in W$. By Schur's lemma
    (Lemma~\ref{lem:starSubalgebra_scalar_commutant_irreducible}) the operator $g$ acts on $W$
    as a scalar, $g x = \lambda\,x$. Pairing with a vector $y \in W$ and using that $P_{W}$ is
    invisible against $W$ gives $\langle f(x), f(y) \rangle = \langle g(x), y \rangle =
    \overline{\lambda}\,\langle x, y \rangle$. Pairing a nonzero $x \in W$ with itself shows
    that $c = \overline{\lambda} = \|f(x)\|^{2} / \|x\|^{2}$, a nonnegative real.
\end{proof}

\begin{lemma}[Nonzero intertwiner scales by a positive factor]%
    \label{lem:starSubalgebra_intertwiner_inner_scale_pos}
    \lean{StarSubalgebra.exists_pos_smul_inner_of_intertwiner_ne_zero}
    \leanok
    \uses{lem:starSubalgebra_intertwiner_inner_scale}
    Let $S$ be a $*$-subalgebra of $\MN{D}$, let $W \subseteq \C^{D}$ be irreducible under $S$,
    and let $f$ be an intertwiner from $W$ into a subspace $W'$ that does not vanish
    on $W$. Then there is a real number $c > 0$ such that for all $x, y \in W$,
    \[
        \langle f(x), f(y) \rangle = c\,\langle x, y \rangle .
    \]
\end{lemma}

\begin{proof}
    \leanok
    \uses{lem:starSubalgebra_intertwiner_inner_scale}
    The nonnegative factor from
    Lemma~\ref{lem:starSubalgebra_intertwiner_inner_scale} satisfies $\|f(x)\|^{2} =
    c\,\|x\|^{2}$. Choosing $x \in W$ with $f(x) \neq 0$ forces both squared norms to be
    positive, so $c > 0$.
\end{proof}

\begin{definition}[Same type for invariant subspaces]%
    \label{def:starSubalgebra_sameIsotype}
    \lean{StarSubalgebra.SameIsotype}
    \leanok
    \uses{def:starSubalgebra_intertwiner}
    Let $S$ be a $*$-subalgebra of $\MN{D}$ and let $W, W' \subseteq \C^{D}$ be invariant
    under $S$. The pieces $W$ and $W'$ are of the same type when there is an intertwiner from
    $W$ into $W'$ that is not zero on $W$.
\end{definition}

\begin{lemma}[Same type is the existence of an isomorphism]%
    \label{lem:starSubalgebra_sameIsotype_iso}
    \lean{StarSubalgebra.sameIsotype_iff_exists_isomorphism}
    \leanok
    \uses{def:starSubalgebra_sameIsotype, lem:starSubalgebra_intertwiner_irreducible}
    Let $S$ be a $*$-subalgebra of $\MN{D}$ and let $W, W' \subseteq \C^{D}$ be irreducible
    under $S$. Then $W$ and $W'$ are of the same type if and only if there is an intertwiner
    from $W$ into $W'$ that is injective on $W$ and maps $W$ onto $W'$.
\end{lemma}

\begin{proof}
    \leanok
    \uses{def:starSubalgebra_sameIsotype, lem:starSubalgebra_intertwiner_irreducible}
    If $W$ and $W'$ are of the same type, then an intertwiner that is nonzero on $W$ is, by
    Schur's lemma (Lemma~\ref{lem:starSubalgebra_intertwiner_irreducible}), injective on $W$
    with image $W'$. Conversely, an intertwiner injective on $W$ with image $W'$ cannot
    vanish on $W$, since $W'$ is nonzero, so the two pieces are of the same type.
\end{proof}

\begin{lemma}[Same-type pieces are unitarily identified]%
    \label{lem:starSubalgebra_sameIsotype_isometry}
    \lean{StarSubalgebra.exists_isometry_intertwiner_of_sameIsotype}
    \leanok
    \uses{lem:starSubalgebra_sameIsotype_iso, lem:starSubalgebra_intertwiner_inner_scale_pos}
    Let $S$ be a $*$-subalgebra of $\MN{D}$ and let $W, W' \subseteq \C^{D}$ be irreducible
    under $S$. If $W$ is of the same type as $W'$, then there is an intertwiner $u$ from $W$
    into $W'$ that maps $W$ onto $W'$ and preserves the inner product on $W$: for all
    $x, y \in W$,
    \[
        \langle u(x), u(y) \rangle = \langle x, y \rangle .
    \]
\end{lemma}

\begin{proof}
    \leanok
    \uses{lem:starSubalgebra_sameIsotype_iso, lem:starSubalgebra_intertwiner_inner_scale_pos}
    Same type gives, by Lemma~\ref{lem:starSubalgebra_sameIsotype_iso}, an intertwiner $f$ that
    is injective on $W$ with image $W'$. Since $W'$ is nonzero, $f$ does not vanish on $W$, so
    by Lemma~\ref{lem:starSubalgebra_intertwiner_inner_scale_pos} there is a real number
    $c > 0$ with $\langle f(x), f(y) \rangle = c\,\langle x, y \rangle$ for all $x, y \in W$.
    Set $u = f / \sqrt{c}$. Scaling by the nonzero factor $1 / \sqrt{c}$ keeps the commuting
    identity and the image $W'$, and turns the inner-product factor into one,
    \[
        \langle u(x), u(y) \rangle = \tfrac{1}{c}\langle f(x), f(y) \rangle = \langle x, y \rangle,
    \]
    so $u$ maps $W$ onto $W'$ and is an isometry on $W$.
\end{proof}

\begin{lemma}[Reflexivity of the same-type relation]%
    \label{lem:starSubalgebra_sameIsotype_refl}
    \lean{StarSubalgebra.sameIsotype_refl}
    \leanok
    \uses{def:starSubalgebra_sameIsotype, def:starSubalgebra_irreducibleSubspace}
    Let $S$ be a $*$-subalgebra of $\MN{D}$ and let $W \subseteq \C^{D}$ be irreducible under
    $S$. Then $W$ is of the same type as itself.
\end{lemma}

\begin{proof}
    \leanok
    \uses{def:starSubalgebra_sameIsotype}
    The identity operator carries $W$ into itself, commutes with every member of $S$, and is
    nonzero on $W$ because $W$ is nonzero.
\end{proof}

\begin{lemma}[Symmetry of the same-type relation]%
    \label{lem:starSubalgebra_sameIsotype_symm}
    \lean{StarSubalgebra.sameIsotype_symm}
    \leanok
    \uses{lem:starSubalgebra_sameIsotype_iso}
    Let $S$ be a $*$-subalgebra of $\MN{D}$ and let $W, W' \subseteq \C^{D}$ be irreducible
    under $S$. If $W$ is of the same type as $W'$, then $W'$ is of the same type as $W$.
\end{lemma}

\begin{proof}
    \leanok
    \uses{lem:starSubalgebra_sameIsotype_iso}
    An isomorphism $f$ of $W$ onto $W'$ has an inverse $g$, obtained by projecting onto $W'$,
    inverting on the subspace, and including back into $W$. For every $A \in S$ and
    $w \in W'$, since $f$ is an intertwiner,
    \[
        g(A w) = f^{-1}(A w) = f^{-1}\bigl(A\,f(f^{-1} w)\bigr)
        = f^{-1}\bigl(f(A\,f^{-1} w)\bigr) = A\,g(w),
    \]
    so $g$ is an intertwiner from $W'$ into $W$ that is injective on $W'$ with image $W$, and
    the two pieces are of the same type in the reverse order.
\end{proof}

\begin{lemma}[Transitivity of the same-type relation]%
    \label{lem:starSubalgebra_sameIsotype_trans}
    \lean{StarSubalgebra.sameIsotype_trans}
    \leanok
    \uses{lem:starSubalgebra_sameIsotype_iso}
    Let $S$ be a $*$-subalgebra of $\MN{D}$ and let $W, W', W'' \subseteq \C^{D}$ be
    irreducible under $S$. If $W$ is of the same type as $W'$ and $W'$ is of the same type as
    $W''$, then $W$ is of the same type as $W''$.
\end{lemma}

\begin{proof}
    \leanok
    \uses{lem:starSubalgebra_sameIsotype_iso}
    Composing an isomorphism $f$ of $W$ onto $W'$ with an isomorphism $g$ of $W'$ onto $W''$
    gives a map $g \circ f$ for which, for every $A \in S$ and $x \in W$,
    \[
        (g \circ f)(A x) = g(f(A x)) = g(A\,f(x)) = A\,g(f(x)) = A\,(g \circ f)(x),
    \]
    so it is an intertwiner from $W$ into $W''$ that is injective on $W$ with image $W''$,
    which is nonzero on $W$ because $W''$ is nonzero.
\end{proof}

\begin{lemma}[The orthogonal projection onto an invariant subspace commutes with the subalgebra]%
    \label{lem:starSubalgebra_projection_commutant}
    \lean{StarSubalgebra.starProjection_toEuclideanLin_comm}
    \leanok
    \uses{lem:starSubalgebra_orthogonal_reducibility}
    Let $S$ be a $*$-subalgebra of $\MN{D}$ and let $W \subseteq \C^{D}$ be invariant under
    $S$, with orthogonal projection $P$ onto $W$. Then for every $A \in S$,
    \[
        P(A x) = A\,(P x)
    \]
    holds for every $x \in \C^{D}$.
\end{lemma}

\begin{proof}
    \leanok
    \uses{lem:starSubalgebra_orthogonal_reducibility}
    Decompose $x = P x + P^{\perp} x$ into its projection onto $W$ and its orthogonal part,
    where $P^{\perp}$ is the projection onto $W^{\perp}$. Each member of $S$ sends $P x$ into
    $W$ (invariance) and $P^{\perp} x$ into $W^{\perp}$ (orthogonal reducibility,
    Lemma~\ref{lem:starSubalgebra_orthogonal_reducibility}), so
    \[
        P(A x) = P\bigl(A\,P x + A\,P^{\perp} x\bigr) = A\,P x + P\bigl(A\,P^{\perp} x\bigr)
        = A\,P x,
    \]
    because $A\,P x \in W$ is fixed by $P$ and $A\,P^{\perp} x \in W^{\perp}$ is annihilated by
    $P$.
\end{proof}

\begin{lemma}[Pieces of different type are orthogonal]%
    \label{lem:starSubalgebra_isOrtho_not_sameIsotype}
    \lean{StarSubalgebra.isOrtho_of_not_sameIsotype}
    \leanok
    \uses{lem:starSubalgebra_intertwiner_irreducible, lem:starSubalgebra_projection_commutant,
        lem:starSubalgebra_sameIsotype_iso}
    Let $S$ be a $*$-subalgebra of $\MN{D}$ and let $W, W' \subseteq \C^{D}$ be irreducible
    under $S$. If $W$ and $W'$ are not of the same type, then $W$ and $W'$ are orthogonal.
\end{lemma}

\begin{proof}
    \leanok
    \uses{lem:starSubalgebra_intertwiner_irreducible, lem:starSubalgebra_projection_commutant,
        lem:starSubalgebra_sameIsotype_iso}
    The orthogonal projection onto $W'$ commutes with every member of $S$
    (Lemma~\ref{lem:starSubalgebra_projection_commutant}) and carries $W$ into $W'$, so it is
    an intertwiner from $W$ into $W'$. By Schur's lemma
    (Lemma~\ref{lem:starSubalgebra_intertwiner_irreducible}) it is either an isomorphism,
    which would make the two pieces of the same type, or zero on $W$. Since the pieces are
    not of the same type, the projection $P_{W'}$ onto $W'$ vanishes on $W$. For every
    $x \in W$ and $y \in W'$ we have $P_{W'} y = y$ and $P_{W'} x = 0$, so self-adjointness
    of $P_{W'}$ gives
    \[
        \langle x, y \rangle = \langle x, P_{W'} y \rangle = \langle P_{W'} x, y \rangle
        = \langle 0, y \rangle = 0,
    \]
    that is, $W$ lies in the orthogonal complement of $W'$.
\end{proof}

Each isotypic component is the supremum of one same-type class; the suprema of two different
classes are orthogonal, and the classes cover the whole family, so the components span
$\C^{D}$.

\begin{lemma}[Suprema of cross-type classes are orthogonal]%
    \label{lem:starSubalgebra_isOrtho_sSup_not_sameIsotype}
    \lean{StarSubalgebra.isOrtho_sSup_of_not_sameIsotype}
    \leanok
    \uses{lem:starSubalgebra_isOrtho_not_sameIsotype}
    Let $S$ be a $*$-subalgebra of $\MN{D}$, and let $\mathcal{D}_1, \mathcal{D}_2$ be two sets
    of subspaces irreducible under $S$. If no piece of $\mathcal{D}_1$ is of the same type as any
    piece of $\mathcal{D}_2$, then
    \[
        \bigvee_{W \in \mathcal{D}_1} W \;\perp\; \bigvee_{W' \in \mathcal{D}_2} W'.
    \]
\end{lemma}

\begin{proof}
    \leanok
    \uses{lem:starSubalgebra_isOrtho_not_sameIsotype}
    Each pair $W \in \mathcal{D}_1$, $W' \in \mathcal{D}_2$ is of different type, hence orthogonal
    by Lemma~\ref{lem:starSubalgebra_isOrtho_not_sameIsotype}. Orthogonality to a fixed subspace
    is preserved by suprema on each side, so the supremum of $\mathcal{D}_1$ is orthogonal to the
    supremum of $\mathcal{D}_2$.
\end{proof}

An irreducible subspace of the supremum of a same-type class need not be one of its pieces,
but its type is determined by the class.

\begin{lemma}[Irreducible subspaces of a supremum inherit a type]%
    \label{lem:starSubalgebra_sameIsotype_of_le_sSup}
    \lean{StarSubalgebra.exists_sameIsotype_of_le_sSup}
    \leanok
    \uses{def:starSubalgebra_irreducibleSubspace, def:starSubalgebra_sameIsotype}
    Let $S$ be a $*$-subalgebra of $\MN{D}$, let $\mathcal{D}$ be a family of subspaces
    irreducible under $S$, and let $W$ be a subspace irreducible under $S$ with
    $W \le \bigvee_{W' \in \mathcal{D}} W'$. Then $W$ is of the same type as some piece of
    $\mathcal{D}$.
\end{lemma}

\begin{proof}
    \leanok
    \uses{lem:starSubalgebra_isOrtho_not_sameIsotype}
    Suppose $W$ is of a different type than every piece of $\mathcal{D}$. Pieces of different
    type are orthogonal (Lemma~\ref{lem:starSubalgebra_isOrtho_not_sameIsotype}), so $W$ is
    orthogonal to every piece of $\mathcal{D}$, hence to their supremum, hence to itself. An
    irreducible subspace is nonzero, a contradiction.
\end{proof}
